\documentclass[12pt,a4paper]{article}
\usepackage{style/main}

\startdocument{Chimie}

\titre{Chapitre 1 - Transformations acido-basiques}{Introduction à la théorie de Brønsted}

Le \textbf{pH} (potentiel hydrogène) mesure la concentration en ions hydronium $\mathrm{H_3O^+}$, souvent notés $\mathrm{H^+}$.  
Dans l'eau, $\mathrm{H^+}$ existe essentiellement sous forme d'ion oxonium $\mathrm{H_3O^+}$.

\section*{I - Théorie de Brønsted}

Selon Brønsted, un \textbf{acide} est une espèce chimique capable de \textbf{céder un proton} $\mathrm{H^+}$, tandis qu'une \textbf{base} est une espèce capable de \textbf{cap­ter un proton} $\mathrm{H^+}$.  
Une réaction acido-basique correspond donc au transfert d'un proton entre l'acide $AH$ d'un couple et la base $B$ d'un autre couple :

\[
AH \rightleftharpoons A^- + H^+
\]

Un \textbf{ampholyte} (ou amphotère) est une espèce capable de jouer le rôle d'acide et de base dans deux couples différents.

\exemple{Acide éthanoïque : $CH_3COOH / CH_3COO^-$}

Un \textbf{amphion} est un ampholyte particulier, qui peut se comporter à la fois comme acide et comme base au sein de la même molécule.  
Par exemple, un acide $\alpha$-aminé :  
\illustration{/home/gladia/Documents/physics-latex/chapitre-1/img/structure_acide_amine.jpg}{Structure générale d'un acide $\alpha$-aminé}
La réaction générale entre l'acide d'un couple et la base d'un autre couple s'écrit :

\[
AH^+ + B^- \rightleftharpoons A + BH
\]

\section*{II - Autoprotolyse de l'eau}

L'eau est un ampholyte et peut réagir avec elle-même (autoprotolyse) :

\[
2 H_2O \rightleftharpoons H_3O^+ + HO^-
\]

Cette réaction n'est pas totale mais elle définit l'équilibre acido-basique de l'eau.

\section*{III - Expressions du pH et produit ionique de l'eau}
\formule{
Le \textbf{pH} d'une solution est défini par :

\[
pH = - \log [H_3O^+]
\]
}
et la concentration en ions hydronium s'écrit :

\[
[H_3O^+] = C_0 \cdot 10^{-pH}
\]

Le produit ionique de l'eau est noté $K_e$ :

\[
K_e = [H_3O^+][HO^-] = 10^{-14} \quad \text{à 25°C}
\]

On définit aussi :

\[
pK_e = - \log K_e = 14
\]
\end{document}